\section{Introduction}
\label{sec:intro}

% --- WR stars as endpoints of massive evolution ---
Wolf-Rayet (WR) stars are evolved, hydrogen-depleted massive stars
characterised by broad emission lines arising from their powerful,
optically thick stellar winds \citep{Crowther2007}. Classical WR stars
(cWR) have shed their hydrogen envelopes, exposing helium-burning cores
with typical masses of $8$--$30$\,\Msun and mass-loss rates of
$\Mdot \sim 10^{-5}$--$10^{-4}$\,\Msun\,yr$^{-1}$. The WR phase
represents one of the final stages before core collapse, making WR stars
direct progenitors of Type~Ibc supernovae and potential progenitors of
stellar-mass black holes \citep{Sander2019}.

% --- WC/WO classification ---
The carbon-sequence WR stars (WC and WO subtypes) are distinguished by
strong carbon and oxygen emission lines, products of core helium burning
brought to the surface by mass loss. They are more chemically evolved than
the nitrogen-sequence (WN) stars, which display CNO-cycle products from
hydrogen burning. The WC/WO phase is therefore a later evolutionary stage,
and these stars are closer to core collapse.

% --- Binary formation channel ---
Two primary channels can produce WR stars: (1)~intrinsic mass loss driven
by radiation pressure on spectral lines, stripping the hydrogen envelope of
sufficiently massive single stars; and (2)~binary interactions, including
Roche-lobe overflow and common-envelope evolution, which strip the envelope
at lower initial masses \citep{Shenar2020}. The relative contribution of
each channel is a fundamental open question in massive star evolution.

% --- Metallicity dependence ---
At solar metallicity, line-driven winds can produce WR stars from single
stars with initial masses $M_\mathrm{ini} \gtrsim 20$\,\Msun. At lower
metallicity, weaker winds raise this threshold ($\gtrsim 25$\,\Msun in the
LMC, $\gtrsim 37$\,\Msun in the SMC), implying that binary interactions
should become increasingly important for WR production
\citep{Shenar2020}. Population synthesis models predict binary fractions
$\gtrsim 50\%$ for WR stars at sub-solar metallicity
\citep{Marchant2023}, but observational constraints remain sparse.

% --- Previous surveys ---
The binary fraction of WR stars has been measured primarily in the Milky
Way. \citet{Dsilva2020} and \citet{Dsilva2023} conducted spectroscopic
multiplicity surveys of Galactic WC/WO and WN stars using
cross-correlation techniques, finding observed binary fractions of
$\sim 40\%$. In the Magellanic Clouds, \citet{Bartzakos2001} surveyed
\Ntotal WC/WO stars in the LMC, detecting only \Nknown spectroscopic
binaries. \citet{Shenar2019} studied the WN stars in the LMC. A
comprehensive survey of the LMC WC/WO population with modern
instrumentation has been lacking.

% --- This work ---
In this paper, we present a spectroscopic multiplicity survey of \Nstars
apparently single WC/WO stars in the LMC, using multi-epoch VLT/X-SHOOTER
observations. We apply the cross-correlation technique of
\citet{Zucker1994} to measure radial velocities and classify binary
candidates. We then perform a Monte-Carlo bias correction following the
methodology of \citet{Dsilva2023} to constrain the intrinsic binary
fraction and orbital period distribution, providing the first such
measurement at LMC metallicity for carbon-sequence WR stars.

The paper is structured as follows. Section~\ref{sec:obs} describes the
sample and observations. Section~\ref{sec:methods} presents the
radial-velocity measurement and binary-classification methodology.
Section~\ref{sec:results} reports the detected binaries and observed binary
fraction. Section~\ref{sec:bias} describes the bias-correction simulation.
Section~\ref{sec:discussion} discusses the results in the context of WR
formation channels and metallicity dependence.
Section~\ref{sec:conclusions} summarises our findings.
